\section{Introduction}\label{sec:intro}

Consider a star and $n\geq2$ planets whose masses are a small multiple $\mu$ of the mass of the star, interacting by Newtonian gravity. To first approximation each planet follows a Kepler ellipse. The mutual attractions deform these ellipses: over times of order $\mu^{-1}$ the eccentricities, inclinations and orientations change appreciably (the secular evolution), while the semi-major axes stay nearly constant. Whether the orbits can come to intersect, so that close encounters, collisions or ejections become possible, depends on how far the eccentricities can grow. In this paper we prove that an elementary conserved quantity bounds this growth for an exponentially long time, with a threshold that is geometrically sharp: at the threshold, two orbits, or an orbit and the star, can touch.

\paragraph{The angular momentum deficit.} The \emph{angular momentum deficit} (AMD) of a planetary system is the difference between the angular momentum that the planets would have on circular, coplanar orbits with the same semi-major axes and their actual angular momentum~\cite[\S2.2]{LP17}. In units where the gravitational constant and the mass of the star are one, and with planet masses $\mu\nu_i$, the Delaunay actions are $\mu\Lambda_i$ with $\Lambda_i=\nu_i\sqrt{a_i}$, and the AMD reads $\mu\sum_i\Lambda_i(1-\sqrt{1-e_i^2}\cos\iota_i)$~\cite[(11)]{LP17}, where $a_i$, $e_i$ and~$\iota_i$ are the semi-major axes, eccentricities and inclinations relative to the plane orthogonal to the total angular momentum. The AMD vanishes exactly for circular coplanar prograde orbits, and a planet can reach large eccentricity only if the AMD is large enough. Laskar and Petit~\cite[\S3]{LP17} computed the least AMD at which the orbits of two planets with given semi-major axes and masses can touch. They call a system \emph{AMD-stable} if its AMD is below this value for every adjacent pair, and \emph{strongly} AMD-stable if, moreover, it is below the circular angular momentum $\mu\Lambda_1$ of the innermost planet, the value at which this planet can reach the star~\cite[\S5 and \S5.3]{LP17}. Since the AMD is conserved by the secular dynamics at all orders, they formulate the following conjecture~\cite[\S3]{LP17}:
\begin{quote}
\emph{we conjecture that in absence of short period resonances, the AMD-stability ensures the practical long-term stability of the system.}
\end{quote}
Here ``practical'' refers to times that are very long compared with the expected lifetime of the star~\cite[\S3, footnote~1]{LP17}. The conjecture is informal: it concerns real systems with actual masses and astronomical time scales, and it is used in~\cite{LP17} to classify observed multi-planet systems.

\paragraph{Chaos and collision safety.} The secular dynamics of a planetary system need not be regular, and the AMD-stability criterion does not ask for its stability. It asks only that the eccentricities and inclinations, however they evolve, never reach a configuration in which orbits touch. In the full Newtonian problem this rests on two facts. First, the total angular momentum~$L$ is conserved exactly, not only in an averaged system. In the normalization of Section~\ref{sec:results}, where the AMD divided by~$\mu$ is $C=\sum_i\Lambda_i-\abs{L}$, this gives
\begin{equation*}
C(t)-C(0)=\sum_{i=1}^n\big(\Lambda_i(t)-\Lambda_i(0)\big)\,,
\end{equation*}
so the AMD can change only as much as the Delaunay actions~$\Lambda_i$, which are functions of the semi-major axes alone. Second, the actions $\Lambda_i$ are conjugate to the fast angles (the mean longitudes), and averaging over these angles shows that they drift very slowly, for a time that is exponentially long in a negative power of~$\mu$. This second fact needs no information about the slow (secular) variables, beyond the requirement that they stay in a region where the Hamiltonian is analytic and the orbits are separated. The AMD, which bounds the eccentricities through $C\geq\sum_i\Lambda_i(1-\sqrt{1-e_i^2})$, keeps the slow variables in such a region. The two facts therefore close a bootstrap: as long as the axes have moved little, the AMD has moved little, so the orbits are still separated and averaging still applies. The secular motion inside the allowed region can be chaotic.

\paragraph{Results.} We make this mechanism rigorous, in the asymptotic regime $\mu\to0$, up to the strong threshold. Let $B(\Lambda)$ be the least of the pair thresholds $B_{i,i+1}(\Lambda)$ and of the actions~$\Lambda_i$, computed exactly from the actions $\Lambda=(\Lambda_1,\ldots,\Lambda_n)$; the inequality $C<B(\Lambda)$ is the strong AMD-stability condition of Laskar and Petit (Section~\ref{sec:results}). Fix the normalized masses $\nu=(\nu_1,\ldots,\nu_n)$, a compact set~$\Aset$ of ordered semi-major axes and a margin $\delta>0$. Informally, our results are as follows.
\begin{itemize}
\item \emph{Theorem~\ref{thm:barrier} (barrier).} For fixed $m>0$, every state whose semi-major axes lie in a compact set of ordered axes and whose AMD is at most $B(\Lambda)-m$ has its perihelia bounded below and adjacent orbital annuli separated by a fixed distance, uniformly. No secular or dynamical information enters.
\item \emph{Theorems~\ref{thm:planar} and~\ref{thm:spatial} (exponential collision safety).} In the plane and in space, if the initial axes lie in~$\Aset$, the initial AMD is at most $B(\Lambda(0))-\delta$, and the initial mean motions satisfy a Diophantine condition with constant $\gamma>0$ and exponent~$\tau>n-1$, then the motion exists and is unique for $\abs{t}\leq T_0\exp(c\mu^{-\nf{1}{(2(\tau+1))}})$. During this time the AMD stays below $B(\Lambda(t))-\nf\delta2$, the axes move by $O(\sqrt\mu)$, and the orbital annuli of the planets remain strictly ordered and uniformly separated from each other and from the star, both for the canonical elements and for the osculating ellipses computed from the actual velocities. No hypothesis on the inclinations is made in space.
\item \emph{Corollary~\ref{cor:half}.} For $\tau=n-\nf12$, which includes almost every initial frequency vector, the time is $T_0\exp(c\mu^{-\nf1{(2n+1)}})$.
\item \emph{Theorem~\ref{thm:resonant} (no arithmetic hypothesis).} Without any condition on the frequencies, the same conclusions hold for $\abs{t}\leq T_0\exp(c\mu^{-\aE{}})$, with drift $O(\mu^{\bE{}})$, where $\aE{}=\nf{2}{(6n^2-3n+14)}$ and $\bE{}=\nf32\,\aE{}$ are Nekhoroshev's exponents for the planetary problem~\cite[(12.4)--(12.5)]{Nek77}.
\end{itemize}
In all results the constants depend only on $\nu$, $\Aset$, $\delta$ (and $\gamma,\tau$), and are chosen before $\mu$ and the initial state. The results are uniform over all eccentricities, inclinations and phases allowed by the AMD condition.

\paragraph{Scope.} These theorems are a rigorous asymptotic form of the mechanism behind the Laskar--Petit conjecture: for every fixed margin below the strong threshold, collision safety holds for exponentially long times once the masses are small enough. They do not resolve the conjecture itself, which concerns practical time scales and actual masses. Our constants are obtained by compactness and are not explicit, and we make no statement about the Solar System or any observed system. We also prove nothing about the stability of the secular dynamics.

\paragraph{Ideas of the proofs.} The barrier (Theorem~\ref{thm:barrier}) follows from the identity for $C$, the inequality $C\geq\sum_i\Lambda_i(1-\sqrt{1-e_i^2})$ and the continuity of~$B$. For Theorems~\ref{thm:planar} and~\ref{thm:spatial} we work in regular Poincar\'e variables, in which the eccentricities and inclinations are Cartesian slow coordinates, and $C$ is half the squared norm of the slow variables. We average over the fast angles about the linear Hamiltonian with the initial frequencies, treating the nonlinear Kepler term on an action ball of radius $\sqrt\mu$ as part of the perturbation. The point of the construction is that the entire set of slow variables allowed by the AMD condition, a ball, is a single domain of analyticity with a fixed complex width. The normal form is of the type of P\"oschel's lemma~\cite[p.~192]{Pos93}; the Diophantine condition with radius $\sqrt\mu$ yields the exponent $\nf{1}{(2(\tau+1))}$, and a first-exit argument closes the bootstrap. Theorem~\ref{thm:resonant} instead applies Nekhoroshev's theorem with slow variables~\cite[Theorem~4.4]{Nek77} on finitely many cells covering the sublevel set $\{C\leq B(\Lambda)-\delta\}$.

\paragraph{Relation to earlier work.} The mechanism is not new, and Theorem~\ref{thm:resonant} is close to a classical result. Nekhoroshev's planetary theorem~\cite[\S12, Theorem~12.3]{Nek77} concerns \emph{domains of planetary motions}: the semi-major axes lie in a box, and the total angular momentum of the unperturbed system, divided by~$\mu$, has norm at least a constant exceeding the largest value for which the boundary of the domain contains collisions of the planets with one another or with the star~\cite[\S12.2, p.~55]{Nek77}. For initial data in such a domain no collision occurs for a time of order $\mu^{-1}\exp(c\mu^{-\aE{}})$, during which the axes change by $O(\mu^{\bE{}})$~\cite[(12.3)--(12.5)]{Nek77}. A lower bound on the angular momentum is an upper bound on the AMD, and the collision value is a strong-AMD threshold taken uniformly over the box. Since $B$ is continuous, covering $\Aset$ by finitely many small boxes appears to give, from Nekhoroshev's theorem, a statement of the type of Theorem~\ref{thm:resonant} under the same hypothesis, up to the margin~$\delta$. What Theorem~\ref{thm:resonant} adds is the identification of the hypothesis with the strong threshold of Laskar and Petit, the conclusions on ordered annuli and osculating ellipses, and a formal proof.

We could not consult Niederman's paper~\cite{Nie96}; our description of it is taken from~\cite{CFG24,BN20}. According to~\cite[\S1.3]{CFG24}, Niederman extended Nekhoroshev theory to the planetary problem, proving that the actions conjugate to the fast angles are stable over exponentially long times provided the secular actions stay in a bounded region, and that near circular coplanar motions the conservation of angular momentum guarantees this. According to~\cite[p.~2]{BN20}, his explicit estimates for the planetary three-body problem near a periodic torus give, at the $5{:}2$ resonance, stability for a time comparable with the age of the Solar System if the ratio of the mass of the larger planet to that of the star is at most $10^{-13}$. Barbieri and Niederman~\cite[\S3.3.2, Theorem~7]{BN20} confine the eccentricities through the conservation of angular momentum in a Nekhoroshev estimate for the planar three-body problem that is local near a periodic orbit of the unperturbed system~\cite[\S3.3.3]{BN20}. Pinzari's construction of Kolmogorov tori in the planetary problem uses an averaging of the fast angles with an exponentially small remainder~\cite[Proposition~5.1]{Pin18}. In the other direction, Clarke, F\'ejoz and Guardia construct orbits of the problem with three planets along which the semi-major axis of the outer planet undergoes large variations~\cite[Theorem~1]{CFG24}; along these orbits the conservation of angular momentum does not confine the secular variables~\cite[\S1.3]{CFG24}.

Theorems~\ref{thm:planar}--\ref{thm:spatial} and Corollary~\ref{cor:half} differ more substantially from these results. Under the Diophantine condition the drift is $O(\sqrt\mu)$ and the time exponent is $\nf{1}{(2(\tau+1))}$, compared with $\mu^{\bE{}}$ and~$\aE{}$. To our knowledge, these statements have not appeared in this sharp form; since we could not consult~\cite{Nie96}, this rests in part on the descriptions above. We do not present the normal form as new.

\paragraph{Formal verification.} Theorems~\ref{thm:barrier}--\ref{thm:resonant} and Corollary~\ref{cor:half} are proved in the proof assistant Lean~4 with the Mathlib library, in the repository~\cite{KAMrepo} at commit \texttt{74d77b43}. The formal statements have the hypotheses, quantifier order and conclusions of Section~\ref{sec:results}; the dynamical ones also list $\Lambda(0)\in\Ach$, which follows from $a(0)\in\Aset$. More detailed proofs are written out in~\cite{KAMnotes}. Appendix~\ref{app:formal} describes what is verified, the correspondence between statements and formal declarations, and what the formal check does not cover.

\paragraph{Organization.} Section~\ref{sec:results} sets up the model and states the results. Section~\ref{sec:barrier} proves the barrier and the identities for the angular momentum and the AMD. Section~\ref{sec:planar} constructs the analytic domain and the normal form and proves Theorem~\ref{thm:planar}, and Section~\ref{sec:spatial} treats the spatial case and Corollary~\ref{cor:half}. Section~\ref{sec:resonant} proves Theorem~\ref{thm:resonant}. Section~\ref{sec:outlook} discusses limitations and open questions, and Appendix~\ref{app:formal} the formalization.
